\documentclass[11pt,a4paper]{article}

\usepackage[T1]{fontenc}
\usepackage[utf8]{inputenc}
\usepackage[ngerman]{babel}
\usepackage{lmodern}
\usepackage{microtype}
\usepackage[a4paper,margin=2.5cm]{geometry}

\usepackage{amsmath,amssymb,amsthm,mathtools}
\usepackage{booktabs}
\usepackage{enumitem}
\usepackage[hidelinks,unicode]{hyperref}

\hypersetup{
  pdftitle={Collatz--Kepler-Gedankenexperiment},
  pdfauthor={Thomas Hoffbauer}
}

\theoremstyle{definition}
\newtheorem{definition}{Definition}[section]
\newtheorem{bemerkung}[definition]{Bemerkung}
\theoremstyle{remark}
\newtheorem{remark}[definition]{Bemerkung}

\newcommand{\Z}{\mathbb{Z}}
\newcommand{\N}{\mathbb{N}}
\newcommand{\Nodd}{\mathbb{N}_{\mathrm{odd}}}
\newcommand{\R}{\mathbb{R}}
\newcommand{\C}{\mathbb{C}}
\newcommand{\Ztwo}{\mathbb{Z}_2}
\newcommand{\Chat}{\widehat{\C}}
\newcommand{\Einf}{E_{\infty}}
\newcommand{\vii}{\nu_2}
\newcommand{\omegaT}{\omega_{12}}
\newcommand{\dist}{\mathrm{dist}}
\newcommand{\SigmaEABC}{\{E,A,B,C\}}
\newcommand{\iotaE}{\iota_{\mathrm{E}}}
\newcommand{\Phpref}{\Phi_{\mathrm{pref}}}

\title{Collatz--Kepler-Gedankenexperiment:\\
  $z_0$, chirale Ellipsen und die Riemann-Kompaktkugel\\
  \large Notizblatt (kein Beweisanspruch)}
\author{Thomas Hoffbauer}
\date{17.\ Juni 2026}

\begin{document}
\maketitle

\begin{abstract}
Dieses Notizblatt verbindet drei Schichten des Collatz- und EABC-Programms:
den komplexen Ankerpunkt $z_0=-1-\mathrm{i}/12$, die chirale
ABCE/CEAB-Dualität der Primvierlingsellipse und die stereographische
Projektion zwischen Riemann-Kompaktkugel und komplexer Ebene.
Zentral ist die wohldefinierte Vorabbildung
$\Phi_{\mathrm{pref}}:\{E,A,B,C\}^{<\omega}\to\C\times\R$ auf
EABC-Wörtern; die dynamische Collatz-Brücke $\Phi:\mathbb{Z}_2\to\mathcal{M}$
folgt erst über eine offene Präfixkodierung $\kappa$.
Es handelt sich um ein \emph{Gedankenexperiment} --- kein Beweisanspruch.
Die Collatz-Vermutung bleibt offen.
\end{abstract}

%======================================================================
\section{Einleitung}
%======================================================================

Die Collatz-Vermutung ist unbewiesen. Bekannte Resultate liefern eine
starke lokale EABC-Grammatik, LTE-Kompensationen und negative
mittlere Driftheuristiken; der punktweise Ausschluss unendlich
schlechter Startwerte ($\Einf=\emptyset$) fehlt weiterhin
(vgl.\ \texttt{collatz\_offene\_punkte.md},
\texttt{collatz\_equivalenz\_e\_infty.tex}).

Dieses Dokument skizziert ein \emph{Gedankenexperiment}: Können
odd-to-odd-Collatz-Bahnen als Kurven in einem Phasenraum
$\mathcal{M}$ gelesen werden, dessen komplexe Koordinate EABC-Phase
und Chiralität trägt und dessen reelle Koordinate Iterationstiefe
bzw.\ $\vii$-Skala misst? Die Konstruktion ist spekulativ und dient
der Suche nach einer konsistenten geometrischen Sprache --- nicht
einem Beweis.

%======================================================================
\section{Drei Ebenen}
%======================================================================

Dieses Notizblatt trennt bewusst drei epistemische Schichten:

\begin{center}
  \begin{tabular}{@{}ll@{}}
    \toprule
    Ebene & Status \\
    \midrule
    Algebraische Identitäten & exakt \\
    Geometrische Interpretation & plausibel \\
    Dynamische Bedeutung & offen \\
    \bottomrule
  \end{tabular}
\end{center}

\noindent
Exakte Aussagen (z.\,B.\ $z_0=\zeta(-1)\cdot(12+\mathrm{i})$,
ABCE/CEAB mod~$12$) dürfen nicht mit geometrischen Lesarten
(Zwölfer-Kreis, Kugelprojektion) oder dynamischen Vermutungen
($\Phpref:\SigmaEABC^{<\omega}\to\mathcal{M}$,
$\Phi:\Ztwo\to\mathcal{M}$, $\mathcal{E}_\infty\subset\mathcal{M}$)
vermischt werden. Zuerst ist die wohldefinierte Wortabbildung
$\Phpref$ zu konstruieren; die Collatz-Brücke $\Phi$ folgt erst
über eine Präfixkodierung $\kappa$ (Abschnitt~\ref{sec:phi}).

%======================================================================
\section{Der Punkt \texorpdfstring{$z_0=-1-\mathrm{i}/12$}{z0=-1-i/12}}
%======================================================================

\subsection{Rolle von \texorpdfstring{$z_0$}{z0}: Anker, nicht Entdeckung}

Der Punkt $z_0$ ist ein \emph{Ankerpunkt} der Projektion: Er fixiert
Orientierung und Skala im komplexen Flavor-Ring. Die eigentliche
offene Aufgabe ist die dynamische Abbildung $\Phi:\Ztwo\to\mathcal{M}$
(Abschnitt~\ref{sec:phi}), vorgebaut als $\Phpref$ auf EABC-Wörtern,
nicht die arithmetische Identität
$z_0=\zeta(-1)\cdot(12+\mathrm{i})$.

\subsection{Anker in der komplexen Ebene}

Setze
\[
z_0 := -1 - \frac{\mathrm{i}}{12}.
\]
Dann gilt
\[
|z_0+1| = \left|-\frac{\mathrm{i}}{12}\right| = \frac{1}{12}.
\]
Der Punkt liegt auf dem Kreis um $-1$ mit Radius $1/12$; er ist
\emph{nicht} der Ursprung und nicht der Nordpol der späteren
Kugelprojektion.

\begin{definition}[Zwölfer-Kreis $K_{12}$]
Der \emph{Zwölfer-Kreis} ist
\[
K_{12}:=\{z\in\C:\ |z+1|=\tfrac{1}{12}\}.
\]
Sein Zentrum ist $-1\equiv C\pmod{12}$; der Radius $1/12$ kodiert
den Nenner der Zwölf-Linie (\texttt{Staudt.tex},
\texttt{zwoelf\_linie\_hoffbauer\_standalone.tex}).
Der Anker $z_0$ liegt auf $K_{12}$ im unteren Halbraum
($\Im z_0=-1/12$).
\end{definition}

\subsection{Verbindung zu \texorpdfstring{$\zeta(-1)$}{zeta(-1)}}

Klassisch ist $\zeta(-1)=-1/12$ (Bernoulli-$B_2$, vgl.\
\texttt{zwoelf\_linie\_hoffbauer\_standalone.tex}). Mit dem
komplexen Nenner $12+\mathrm{i}$ folgt die Identität
\[
z_0 = \zeta(-1)\,(12+\mathrm{i}) = -\frac{1}{12}\,(12+\mathrm{i})
    = -1 - \frac{\mathrm{i}}{12}.
\]
Der Nenner $12$ kodiert den mod-$12$-Flavor-Ring; die imaginäre
Komponente $\mathrm{i}$ verweist auf die quaternionische Phase der
EABC-Kante $E\leftrightarrow C$ (\texttt{energiedoku\_eabc\_c4\_kohaerenz.tex}).

\subsection{Restklasse und $C$-Typ}

Es gilt $\Re(z_0)=-1\equiv 11\pmod{12}$. Im EABC-Modell entspricht
dies der Klasse $C$ ($C\equiv 11\pmod{12}$). Der Anker $z_0$ liegt
damit auf der \emph{chiralen Startseite CEAB} des Flavor-Rings
($p\equiv 11\Rightarrow\mathrm{CEAB}$, Satz in
\texttt{Projektionszeuge\_Primvierling.tex}).

\subsection{\texorpdfstring{$C$}{C}-Klasse und chirale Asymmetrie}
\label{sec:c-asymmetrie}

Die $C$-Klasse ($\equiv 11\pmod{12}$) ist die CEAB-Startseite des
Flavor-Rings. Zwei Asymmetrien sind zu unterscheiden:

\paragraph{Bewiesen (Ebene~A).}
Markierte Primvierlinge mit $p\equiv 5$ tragen $\mathrm{ABCE}$,
mit $p\equiv 11$ tragen $\mathrm{CEAB}$; CEAB ist die zyklische
Verschiebung von ABCE (\texttt{Projektionszeuge\_Primvierling.tex},
\texttt{EABC.lean}).

\paragraph{Geometrisch plausibel (Ebene~B).}
In der \#Energiedoku gilt
$\overline{\mathrm{ABCE}}=\mathrm{CEAB}$ als komplex konjugierte
Phasen: $\chi_H(\mathrm{ABCE})=+\mathrm{i}$,
$\chi_H(\mathrm{CEAB})=-\mathrm{i}$
(\texttt{energiedoku\_eabc\_c4\_kohaerenz.tex}).
Die beiden Ellipsen-Lesarten sind chiral konjugiert mit gleicher
Exzentrizität $e_{\mathrm{PV}}=\sqrt{3}/2$.

\paragraph{Offen (Ebene~C).}
Globale Zählasymmetrie $\#\mathrm{ABCE}(N)$ vs.\ $\#\mathrm{CEAB}(N)$:
Siebstatistiken bis $10^8$ zeigen keinen robusten Bias
($z\lesssim 0{,}7$, \texttt{Projektionszeuge\_Primvierling.tex},
Abschnitt~Ebene~C). Der Anker $z_0$ auf der $C$-Seite des
Zwölfer-Kreises ist damit \emph{kein} Beweis einer $C$-Dominanz,
sondern eine Orientierungswahl.

\subsection{2-adische Lesart}

In $\Ztwo$ ist $-1$ der unendliche Bitstrom $\ldots 111_2$.
Collatz-Bahnen werden 2-adisch als Präfixe in $\Ztwo$ studiert
(\texttt{collatz\_z2\_attraktor.tex}). Die Gleichung $z_0=\zeta(-1)\cdot(12+\mathrm{i})$
verknüpft den rationalen Bernoulli-Anker $-1/12$ mit der
arithmetischen $-1$-Restklasse modulo $12$; eine dynamische
Kopplung zu $\dist_2(\cdot,\Einf)$ ist hier \emph{nicht} bewiesen,
nur als Lesart notiert.

%======================================================================
\section{Zwei chiral konjugierte Ellipsen (ABCE / CEAB)}
%======================================================================

\subsection{Primvierlingsellipse}

Für einen Primvierling $Q(p)=(p,p+2,p+6,p+8)$ mit Symmetrieanker
$M=p+4$ definiert \texttt{Projektionszeuge\_Primvierling.tex} die
Bamberger Primvierlingsellipse mit
\[
a_{\mathrm{PV}}=4,\qquad b_{\mathrm{PV}}=2,\qquad
f_{\mathrm{PV}}=2\sqrt{3},\qquad
e_{\mathrm{PV}}=\frac{\sqrt{3}}{2}.
\]
Die vier Primpositionen bilden das symmetrische Fenster
$\{M-4,M-2,M+2,M+4\}$; die Schrittfolge $2$-$4$-$2$ ist die
arithmetische Kepler-Ellipse des Vierlings.

\subsection{Chirale Konjugation}

Markierte Primvierlinge tragen genau zwei Wörter:
\[
p\equiv 5\pmod{12}\Rightarrow \mathrm{ABCE},\qquad
p\equiv 11\pmod{12}\Rightarrow \mathrm{CEAB}.
\]
$\mathrm{CEAB}$ ist die zyklische Verschiebung von $\mathrm{ABCE}$
entlang $E\to A\to B\to C\to E$; die Orientierung ist eine
Phasendifferenz um $\pi$ auf dem viergliedrigen Flavor-Ring.

In der \#Energiedoku (\texttt{energiedoku\_eabc\_c4\_kohaerenz.tex})
werden die Hurwitz-Phasen
\[
\chi_H(\mathrm{ABCE})=+\mathrm{i},\qquad
\chi_H(\mathrm{CEAB})=-\mathrm{i}
\]
numerisch verifiziert. Die beiden Ellipsen-Lesarten sind damit
\emph{chiral konjugiert}: gleiche Exzentrizität $e_{\mathrm{PV}}=\sqrt{3}/2$,
entgegengesetzte quaternionische Phase ($\pm\mathrm{i}$).

\begin{bemerkung}
Der Wert $3/2$ in Ebene~B des Projektionszeugen ist das
projektive Kantengewicht $\rho_{\mathrm{PV}}$, nicht die
Ellipsenexzentrizität ($e_{\mathrm{PV}}<1$).
\end{bemerkung}

%======================================================================
\section{Riemann-Kompaktkugel und komplexe Ebene}
%======================================================================

\subsection{Stereographische Projektion}

Identifiziere die Einheitskugel
\[
S^2=\{(x,y,z)\in\R^3:\ x^2+y^2+z^2=1\}
\]
mit der Riemann-Kompaktkugel $\Chat=\C\cup\{\infty\}$.
Mit Nordpol $N=(0,0,1)$, Äquator $\{z=0\}$ und Südpol $S=(0,0,-1)$
definiere die stereographische Projektion
\[
\sigma:\ S^2\setminus\{N\}\longrightarrow \C,\qquad
\sigma(x,y,z)=\frac{x+\mathrm{i}y}{1-z},
\]
sowie $\sigma(N)=\infty$. Die Umkehrabbildung für $w\in\C$ ist
\[
\iota(w)=\left(
  \frac{2\Re w}{1+|w|^2},\
  \frac{2\Im w}{1+|w|^2},\
  \frac{|w|^2-1}{1+|w|^2}
\right)\in S^2\setminus\{N\}.
\]
Diese Konvention entspricht \texttt{Edinburg.py} und der
Mathlib-Spur \texttt{Geometry.Manifold.Instances.Sphere}
(vgl.\ \texttt{collatz\_mathlib\_eabc\_kandidaten.md}).

\subsection{Bezug zur Hurwitz-Doppelkugel}

Im Forschungsprogramm \texttt{document.tex} operieren Primvierlinge
dual auf den Schalen der Hurwitz-Doppelkugel:
\[
\Phi_T:\ \mathbb{P}_4\to \mathrm{Aut}(\mathbb{O})\cong G_2
\quad\text{(Torsion, integrale Schale $\mathcal{S}_{\mathbb{Z}}$)},
\]
\[
\Phi_R:\ \mathbb{P}_4\to \mathrm{Spin}(8)
\quad\text{(Translation, halbintegrale Schale $\mathcal{S}_{\mathbb{Z}+1/2}$)}.
\]
Die Kompaktkugel $S^2$ ist hier die \emph{projektive Schattenebene}
der Bamberg-Kugel: Richtungen in $\mathbb{H}$ bzw.\ $\mathbb{O}$
werden stereographisch auf $\C$ abgebildet, während EABC-Phasen die
$G_2$- bzw.\ Spin$(8)$-Umläufe kodieren
(\texttt{Grundsatzartikel\_Hurwitz\_Raum.tex}, \texttt{kapitel-eabc-phaenomenologie.tex}).
Eine kausale Verknüpfung $\Phi_T,\Phi_R\leftrightarrow\sigma$ ist
offen; dieses Notizblatt nutzt $\sigma$ nur als explizite
Koordinatenwahl.

\subsection{Lage von \texorpdfstring{$z_0$}{z0} auf \texorpdfstring{$S^2$}{S2}}

Für $z_0=-1-\mathrm{i}/12$ mit $|z_0|^2=145/144$ ergibt die Umkehrung
\[
P_0:=\iota(z_0)=\left(
  -\frac{288}{289},\
  -\frac{24}{289},\
  \frac{1}{289}
\right)\in S^2.
\]
Der Punkt liegt im nördlichen Halbraum ($z>0$), nahe am Äquator
($z=1/289\ll 1$), im dritten Quadranten der $(x,y)$-Ebene.
Er ist weder Nord- noch Südpol; unter $\sigma$ entspricht er genau
$z_0\in\C$. Die Nähe zum Südpol $S$ (Bild $0$) spiegelt die
Kleinheit von $|z_0+1|=1/12$.

%======================================================================
\section{Phasenraum \texorpdfstring{$\mathcal{M}=\C\times\R$}{M=C×R}}
%======================================================================

\begin{definition}[Basisfasenraum]
Setze $\mathcal{M}:=\C\times\R$. Die Koordinate $z\in\C$ trägt
EABC-Phase und Chiralität ($\omegaT^{k}$ mit
$\omegaT=e^{2\pi\mathrm{i}/12}$, vgl.\ \texttt{BM Messprotokoll.tex});
die Koordinate $t\in\R$ misst Iteration, logarithmische Tiefe oder
$\vii$-akkumulierte 2-adische Distanz entlang einer Collatz-Bahn.
\end{definition}

Eine odd-to-odd-Bahn $(n_k)_{k\geq 0}$ mit $n_{k+1}=U(n_k)$ lässt sich
\emph{heuristisch} als Kurve
\[
\gamma(k)=\bigl(z_k,\ t_k\bigr)\in\mathcal{M}
\]
lesen, wobei $z_k$ aus dem EABC-Wort $w(n_k)\in\{E,A,B,C\}^{\N}$
via $z_k\sim\omegaT^{\mathrm{index}(w_k)}$ konstruiert wird und
$t_k$ z.\,B.\ $t_k=\log n_k$ oder $t_k=\sum_{j<k}\vii(3n_j+1)$ wählt.

\subsection{Erweiterung über die Kugel}

Alternativ arbeite auf
\[
\widetilde{\mathcal{M}}:=S^2\times\R
\quad\text{oder}\quad
\Chat\times\R
\]
mit Projektion $\pi_{\C}:\widetilde{\mathcal{M}}\to\mathcal{M}$,
$\pi_{\C}(\iota(w),t)=(w,t)$. Der Anker $P_0\in S^2$ fixiert eine
\emph{globale Orientierung} der Kugel; lokale Collatz-Schritte
ändern $t$ und rotieren $z$ auf dem Flavor-Ring, ohne den Pol zu
berühren (für endliche Bahnen mit $|z+1|\gg 0$).

%======================================================================
\section{Wohldefinierte Vorabbildung \texorpdfstring{$\Phpref$}{Phi\_pref}}
\label{sec:phi}
%======================================================================

\subsection{EABC-Alphabet und Restklassen modulo~$12$}

Das EABC-Alphabet ist $\SigmaEABC=\{E,A,B,C\}$.
Jeder Buchstabe kodiert eine primfähige Restklasse modulo~$12$
(\texttt{EABC.lean}, \texttt{CollatzEabc.Mod12Matrix}):
\[
E\equiv 1,\quad A\equiv 5,\quad B\equiv 7,\quad C\equiv 11\pmod{12}.
\]
Endliche Wörter liegen in $\SigmaEABC^{<\omega}=\bigcup_{n\in\N}\SigmaEABC^n$;
$|w|$ bezeichnet die Wortlänge.

\subsection{Phasenindex, \texorpdfstring{$\Theta(w)$}{Theta(w)} und Einheitsphase \texorpdfstring{$u(w)$}{u(w)}}

\begin{definition}[Phasenindex und kumulative Phase]
Der \emph{Phasenindex} ist
\[
\iotaE:\ \SigmaEABC\longrightarrow \{0,1,2,3\},\qquad
E\mapsto 0,\ A\mapsto 1,\ B\mapsto 2,\ C\mapsto 3.
\]
Für $w=w_0\cdots w_{n-1}\in\SigmaEABC^n$ setze
\[
\Theta(w):=\sum_{j=0}^{n-1}\iotaE(w_j)\ \in\ \Z,
\qquad
\Theta_{4}(w):=\Theta(w)\bmod 4.
\]
Die \emph{Einheitsphase} ist
\[
u(w):=\exp\!\Bigl(\mathrm{i}\,\frac{\pi}{2}\,\Theta_{4}(w)\Bigr)
\ \in\ \{1,\mathrm{i},-1,-\mathrm{i}\}.
\]
\end{definition}

\begin{definition}[Chiralitätswörter]
Die kanonischen Chiralitätswörter sind
$\mathrm{ABCE}$ und $\mathrm{CEAB}$ (Abschnitt~\ref{sec:c-asymmetrie},
\texttt{Projektionszeuge\_Primvierling.tex}).
\end{definition}

\begin{bemerkung}[Chiralitätsaxiom]
Die Hurwitz-Phasen erfüllen
$\chi_H(\mathrm{ABCE})=+\mathrm{i}$ und
$\chi_H(\mathrm{CEAB})=-\mathrm{i}$
(\texttt{energiedoku\_eabc\_c4\_kohaerenz.tex}).
Wir fixieren die Vorabbildung so, dass für die Chiralitätswörter
explizit gilt:
\[
u(\mathrm{ABCE})=+\mathrm{i},\qquad
u(\mathrm{CEAB})=-\mathrm{i}=-\,u(\mathrm{ABCE}).
\]
Für allgemeine Wörter wird $u(w)$ über $\Theta_{4}$ und den
Startbuchstaben gelesen; die Konsistenz mit $\chi_H$ auf
$\{\mathrm{ABCE},\mathrm{CEAB}\}$ ist die chirale Normalisierung.
\end{bemerkung}

\subsection{Anker \texorpdfstring{$z_0$}{z0} und Zwölfer-Kreis \texorpdfstring{$K_{12}$}{K12}}

Der Anker $z_0=-1-\mathrm{i}/12$ und der Zwölfer-Kreis
$K_{12}=\{z\in\C:\ |z+1|=1/12\}$ sind in
Abschnitt~3 definiert. Sie fixieren Orientierung und radiale
Skala der Vorabbildung; $z_0\in K_{12}$ liegt auf der
$C$-Seite (CEAB-Start, $C\equiv 11\pmod{12}$).

\subsection{Radialkoordinate und komplexe Spur}

Für $w\in\SigmaEABC^{<\omega}$ setze die \emph{Radialkomponente}
\[
\rho(w):=2^{-|w|}\in(0,1],
\qquad
z(w):=z_0+\rho(w)\,u(w)\in\C.
\]
Der Punkt $z(w)$ liegt in einer $\rho(w)$-Umgebung von $z_0$;
für $|w|\to\infty$ konvergiert $z(w)\to z_0$ auf $K_{12}$.

\subsection{Zeitkoordinate und Vorabbildung \texorpdfstring{$\Phpref$}{Phi\_pref}}

\begin{definition}[Vorabbildung auf EABC-Wörtern]
Die reelle \emph{Zeitkoordinate} ist $t(w):=|w|\in\N$.
Die \emph{Vorabbildung} ist
\[
\boxed{\;
\Phpref:\ \SigmaEABC^{<\omega}\longrightarrow \mathcal{M}=\C\times\R,\qquad
\Phpref(w)=\bigl(z(w),\,t(w)\bigr).
\;}
\]
\end{definition}

\noindent
Damit ist $\Phpref$ auf rein kombinatorischen EABC-Wörtern
\emph{wohldefiniert} --- ohne Collatz-Dynamik, ohne $\Ztwo$-Einbettung
und ohne Behauptung über Bahnen.

\subsection{Einschränkung: Wörter zuerst, Collatz erst über Kodierung}

$\Phpref$ lebt auf $\SigmaEABC^{<\omega}$, nicht auf $\Ztwo$.
Eine Collatz-Anbindung erfordert eine \emph{Präfixkodierung}
\[
\kappa:\ \Ztwo\longrightarrow \SigmaEABC^{<\omega}
\quad\text{(bzw.\ }\kappa_K:\Ztwo\to\SigmaEABC^K\text{ für festes }K\text{)},
\]
sodass
\[
\Phi_K:=\Phpref\circ\kappa_K:\ \Ztwo\longrightarrow \mathcal{M}.
\]
Die Konstruktion von $\kappa$ (aus odd-to-odd-Bahnen, 2-adischen
Präfixen oder endlichen Grammatikfenstern) ist \emph{offen} und
getrennt von der Definition von $\Phpref$.
Spekulative Kandidaten ($\vii$-Summen, LTE-Blöcke,
$\|x\|_2^{-1}$) bleiben Heuristiken für $\kappa$, nicht für
$\Phpref$ selbst.

\subsection{Lean-kompatibler Kern}

Statt $\exp(\mathrm{i}\theta)$ kodieren die vier Einheitsphasen
als Gitterpaare in $\Z\times\Z$ (\texttt{CollatzEabc.PrefProjection}):
\[
1\mapsto(1,0),\quad
\mathrm{i}\mapsto(0,1),\quad
-1\mapsto(-1,0),\quad
-\mathrm{i}\mapsto(0,-1).
\]
Die Abbildung $\mathrm{phaseUnit}:\Z/4\Z\to\Z\times\Z$ ist
sorry-frei formalisierbar; die komplexe Spur $z(w)$ bleibt
im TeX-Notizblatt, bis eine $\C$-Einbettung in Mathlib
explizit gewünscht ist.

\subsection{Brücke zu \texorpdfstring{$\Phi:\Ztwo\to\mathcal{M}$}{Phi: Z2→M} und \texorpdfstring{$\Einf$}{E\_infty}}

Die zentrale \emph{dynamische} offene Aufgabe bleibt
\[
\Phi:\ \Ztwo\longrightarrow \mathcal{M},
\qquad
\Phi=\Phpref\circ\kappa
\]
für eine noch zu konstruierende Kodierung $\kappa$.
Collatz ist äquivalent zu
\[
\Einf=\emptyset
\quad\Longleftrightarrow\quad
\text{keine ungerade Bahn erreicht den trivialen Fixpunkt }1,
\]
spekulativ gelesen als: jede Bahn $\gamma(k)=\Phi(U^k(n))$ in
$\mathcal{M}$ trifft irgendwann die Fixpunktregion, statt mit
$t\to\infty$ in $\mathcal{E}_\infty$ zu verharren
(vgl.\ \texttt{collatz\_equivalenz\_e\_infty.tex}).

\subsection{\texorpdfstring{$\Einf$}{E\_infty} als geometrisches Objekt in \texorpdfstring{$\mathcal{M}$}{M}}

Die echte Ausnahmemenge
\[
\Einf=\{n\in\Nodd:\forall K,\ U^K(n)\neq 1\}
\]
(\texttt{collatz\_equivalenz\_e\_infty.tex}) ist Collatz-äquivalent zu
$\Einf=\emptyset$. Spekulativ definiere das geometrische Bild
\[
\mathcal{E}_\infty:=\overline{\Phi(\Einf)}\subset\mathcal{M}
\]
(bzw.\ $\overline{\Phi(\Einf)}\subset S^2\times\R$).
\paragraph{Lesart.}
\begin{itemize}[nosep]
  \item $\Einf=\emptyset$: $\mathcal{E}_\infty=\emptyset$; alle Bahnen
    verlassen die invariante Schleife und erreichen die Fixpunktregion.
  \item Hypothetisches Gegenbeispiel: $\mathcal{E}_\infty$ wäre eine
    unbeschränkte Kurven- oder Flächenmenge mit $t\to\infty$ und
    asymptotisch festem Phasenband auf $K_{12}$ oder einem
    ABCE-invarianten Torus in $\mathcal{M}$.
\end{itemize}
Keine Aussage dieser Art ist bewiesen; $\mathcal{E}_\infty$ ist ein
\emph{geometrisches Platzhalterobjekt} für die offene Collatz-Brücke.

%======================================================================
\section{Bernoulli-Uhr auf einer chiralen Doppelellipse}
\label{sec:bernoulli-uhr}
%======================================================================

\begin{definition}[Bernoulli-Uhr]
Sei $z_0=-1-\mathrm{i}/12$, $r_m=2^{-m}$, $\theta_m=\pi/2\cdot m$.
Für $m\geq 1$:
\[
U(m)=\bigl(
z_0-r_m B_{2m-2}e^{\mathrm{i}\theta_m},\;
z_0+r_m B_{2m}e^{\mathrm{i}\theta_m},\;
z_0+r_m B_{2m+2}e^{\mathrm{i}\theta_m}
\bigr),
\qquad
B_{2m}=-2m\,\zeta(1-2m).
\]
Dreierschritt: $3\cdot 2\pi/12=\pi/2$.
\end{definition}

\begin{remark}[Chirale Interpretation]
Vierteltakt $E\to A\to B\to C\to E$, ABCE vs.\ CEAB um $\pi$,
$z_0=\zeta(-1)(12+\mathrm{i})$, geometrische Organisation $\{B_{2m}\}$.
\end{remark}

\begin{remark}[Nichtbehauptung]
\begin{itemize}[nosep]
  \item nicht Collatz,
  \item nicht Konjugation,
  \item nicht $E_\infty$,
  \item nicht Beweis.
\end{itemize}
\end{remark}

\subsection{arXiv- und Lean-Belege (ehrlich)}

\paragraph{Unterstützt (algebraisch / 2-adisch).}
\begin{itemize}[nosep]
  \item Bernoulli--Zeta: arXiv:math/0411498 (Struktur der Bernoulli-Zahlen,
    Kubota--Leopoldt); klassisch von Staudt--Clausen.
  \item Stirling--Bernoulli--2-adisch: arXiv:1912.01090, arXiv:2111.08766
    ($\nu_2$ von Stirling-Zahlen via Bernoulli-Polynome).
  \item 2-adische Collatz-Parität: arXiv:2601.12772 (Ghost Cycles, Paritätsmuster).
  \item Lean/Mathlib: \texttt{NumberTheory.Bernoulli} (\texttt{bernoulli},
    \texttt{bernoulli\_two}), \texttt{PrefProjection.lean}
    (\texttt{Fin~4}, \texttt{List.fold}/\texttt{sum}, \texttt{rho}$=2^{-|w|}$).
\end{itemize}

\paragraph{Spekulativ / nicht belegt.}
\begin{itemize}[nosep]
  \item Chirale Bernoulli-Uhr auf Ellipsenachsen $\mathcal{E}_\pm$.
  \item Dreiertakt $\Leftrightarrow$ Collatz-Dreierschritt.
  \item Hurwitz-Quaternionen mod~$12$ als Uhrwerk
    (arXiv:1201.5817, arXiv:0806.3943: Faktorisierung, keine mod-$12$-Uhr).
  \item $\zeta(1-2m)$-Nenner~$12$ $\Rightarrow$ dynamische Collatz-Konvergenz.
\end{itemize}

%======================================================================
\section{Stirling, Bernoulli und triviale Nullstellen}
\label{sec:stirling-bernoulli}
%======================================================================

\subsection{Exakte algebraische Verknüpfungen (im Repo dokumentiert)}

\paragraph{$\zeta(-1)$ und Bernoulli.}
Klassisch $\zeta(-1)=-B_2/2=-1/12$
(\texttt{Staudt.tex}, \texttt{zwoelf\_linie\_hoffbauer\_standalone.tex}).
Allgemein für ungerade negative Stellen
\[
\zeta(-(2m-1))=-\frac{B_{2m}}{2m}.
\]

\paragraph{Triviale Nullstellen.}
Für gerade negative $n=2,4,6,\ldots$ gilt $\zeta(-n)=0$; die
\emph{trivialen Nullstellen} liegen äquidistant auf der negativen
reellen Achse bei $-2,-4,-6,\ldots$ (Standardanalysis; im Repo
nicht als Collatz-Objekt weiterverarbeitet).

\paragraph{Negative gerade Spezialwerte via Bernoulli.}
Für positive gerade $2n$:
$\zeta(2n)=(-1)^{n+1}B_{2n}(2\pi)^{2n}/(2(2n)!)$; die Bernoulli-Zahlen
$|B_{2n}|$ wachsen asymptotisch via Stirling
(\texttt{collatz\_lean\_mathlib\_update.tex}: Robbins-Schranke
\texttt{Stirling.log\_stirlingSeq\_diff\_le}, Kombination $|B_{2n}|$
via Stirling noch offen in Mathlib).

\paragraph{Bernoulli-Normschalen (No-Go).}
Via von Staudt--Clausen:
$\mathcal{I}(n):=|B_{2n}|\cdot\prod_{(p-1)\mid 2n}p$
(\texttt{collatz\_schlussartikel\_arxiv.tex}, §Bernoulli-Normschalen).
Unter $n\mapsto 3n+1$ wächst $\mathcal{I}$ super-exponentiell ---
\emph{kein} Lyapunov-Kandidat für Collatz.

\paragraph{Hoffbauers Zwölf-Linie.}
Die Menge $\mathcal{Z}_{12}=\{m:Q_m=12\}$ der Zwölf-Indizes
(\texttt{Staudt.tex}) markiert ungerade negative Zetawerte mit
Minimalnenner $12$; $\zeta(-1)=-1/12$ ist die \emph{Urform}.
Die Bernoulli-``Brüder'' höherer Ordnung ($B_4,B_6,\ldots$) erzeugen
über $-B_{2m}/(2m)$ die rationalen Zetawerte auf ungeraden negativen
Stellen; ihre Nennerstruktur wird von von Staudt--Clausen gesteuert.

\subsection{Stirling als Bulk-Hintergrund}

In \texttt{Sterling.tex} ist Stirling der glatte Bulk-Hintergrund
($\log(n!)\sim n\log n-n+\cdots$), EABC der diskrete Restklassenkanal.
Die Verbindung Bernoulli$\leftrightarrow$Stirling ist in der
Literatur standard (Asymptotik $|B_{2n}|\sim 4\sqrt{\pi n}
(n/\pi e)^{2n}$); im Repo als Mathlib-Baustein notiert, aber nicht
als Collatz-Mechanismus formalisiert.

\subsection{Äquidistanz: ehrliche Abgrenzung}

\begin{center}
  \begin{tabular}{@{}lll@{}}
    \toprule
    Objekt & Äquidistanz? & Status im Repo \\
    \midrule
    Triviale Nullstellen $-2,-4,-6,\ldots$ & ja (Abstand $2$) & Standard, nicht Collatz-verknüpft \\
    Hoffbauers Zwölf-Linie $\mathcal{Z}_{12}$ & nein (arithmetisch dünn) & \texttt{Staudt.tex}, exakt charakterisiert \\
    Nichttriviale Nullstellen (krit.\ Linie) & GUE-Statistik, nicht exakt äquidistant & \texttt{Resonator.tex}, \texttt{rh scholze.tex} \\
    Punkte auf $K_{12}$ & nein (ein Kreis, kein Gitter) & Spekulation dieses Notizblatts \\
    \bottomrule
  \end{tabular}
\end{center}

\noindent
Eine Verknüpfung ``äquidistante triviale Nullstellen $\leftrightarrow$
Bernoulli-Brüder $\leftrightarrow$ Zwölfer-Kreis'' ist im Repo
\emph{nicht} dokumentiert. Die Analogie ist heuristisch: Die trivialen
Nullstellen bilden ein äquidistantes Gitter auf $\Re(s)<0$; die
Bernoulli-Zahlen kodieren die \emph{rationalen} Werte auf den
dazwischenliegenden ungeraden negativen Stellen; $K_{12}$ ist eine
geometrische Lesart des Nenners $12$ bei $\zeta(-1)$, nicht ein
Nullstellenort der Zetafunktion.

%======================================================================
\section{Numerische Zeugen und Nicht-Zeugen}
\label{sec:numerische-zeugen}
%======================================================================

\begin{definition}[Epistemische Taxonomie numerischer Zeugen]
Im EABC-Modell heißt ein \emph{Zeuge} eine reproduzierbare numerische
oder geometrische Beobachtung, die \emph{nur} die EABC-Struktur
(Restklassen mod~$12$, Primvierlingsfenster, Projektionsoperator,
polyedrische Symmetrie) stützt --- ohne Collatz-Dynamik,
ohne $\Einf$ und ohne Bernoulli-Lyapunov-Anspruch.
\end{definition}

\begin{bemerkung}[Abgrenzung]
Die folgende Tabelle trennt bewusst \emph{was belegt} von
\emph{was nicht belegt}. Ein Eintrag unter ``Zeuge JA'' ist kein
Schritt zum Collatz-Beweis; ein Eintrag unter ``Zeuge NEIN'' kann
wohldefiniert, aber ohne numerischen Befund oder nur Gedankenmodell sein.
\end{bemerkung}

\begin{center}
  \small
  \begin{tabular}{@{}p{2.6cm}p{1.4cm}p{4.8cm}p{4.8cm}@{}}
    \toprule
    Zeuge & Status & Was belegt & Was \emph{nicht} belegt \\
    \midrule
    Chiralität ABCE/CEAB &
    JA &
    Chirale Doppelstruktur markierter Primvierlinge; im Integrationsstrom
    $\sim\!100$ vs.\ $\sim\!26$ gegen Permutations-Null ($p<10^{-4}$);
    relative Signatur $\approx 47\%/53\%$ mit Null vereinbar
    (\texttt{wolfram.py}, \texttt{SCIENTIFIC\_CLAIM}) &
    Collatz; globale $C$-Dominanz; dynamische $\Phi$-Brücke \\
    \addlinespace
    Span-Zeuge &
    JA &
    $\mathrm{span}(Q(p))=p_4-p_1=10+12g$ für Primvierlingsfenster;
    modulare Leiter-Geometrie (\texttt{Beweis\_Spanne\_EABC.tex},
    \texttt{prognoseEABC.py}) &
    Punktweise Collatz-Konvergenz; $\Einf=\emptyset$ \\
    \addlinespace
  $24I_3$-Zeuge &
    JA &
    $M_{\mathrm{eff}}\to 24I_3$ als isotroper Fixpunkt unter
    $R^*_{\mathrm{EABC}}$ (\texttt{EABC Intermediate.tex},
    \texttt{eabc\_icosahedron\_test.py}) &
    Collatz-Attraktor $E$; Bernoulli-Uhr als Dynamik \\
    \addlinespace
    Ikosaeder/Dodekaeder &
    JA &
    Anisotropy $=0$ für geometrische EABC-Zuweisung vs.\ $>0$ bei
    Zufallszuweisung (\texttt{eabc\_icosahedron\_test.py}) &
    Hurwitz-Quaternionen-Uhr; mod-$12$-Spezialfall allein \\
    \midrule
    Bernoulli-Uhr &
    NEIN &
    Wohldefinierte Gewichtsformel $U(m)$ auf $\mathcal{E}_\pm$
    (Abschnitt~\ref{sec:bernoulli-uhr}) &
    Numerischer Collatz-Befund; Konjugation; $\Einf$ \\
    \addlinespace
    $z_0$ dynamisch &
    NEIN &
    Algebraische Identität $z_0=\zeta(-1)(12+\mathrm{i})$,
    Anker auf $K_{12}$ (Abschnitt~3) &
    Relevanz für $\dist_2(\cdot,\Einf)$ oder $\Phi$ \\
    \addlinespace
    $\Phpref$ &
    offen &
    Wohldefinierte Vorabbildung auf $\SigmaEABC^{<\omega}$
    (Abschnitt~\ref{sec:phi}, \texttt{PrefProjection.lean}) &
    Dynamische Brücke $\Phi=\Phpref\circ\kappa$;
    \emph{Datentest} unten \\
    \addlinespace
    Collatz / $\Einf$ &
    NEIN &
    Präzise Äquivalenz $\Einf=\emptyset\Leftrightarrow$ Collatz
    (\texttt{collatz\_equivalenz\_e\_infty.tex}) &
    Beweis; Uniformitätslücke (§2 in
    \texttt{collatz\_offene\_punkte.md}) \\
    \bottomrule
  \end{tabular}
\end{center}

\subsection{Offener Datentest: \texorpdfstring{$\Phpref$}{Phi\_pref}}

\begin{definition}[Datentest $\Phpref$ auf Primvierlingswörtern]
Für $w\in\SigmaEABC^{<\omega}$ mit Phasenindex $\Theta(w)$ und
chiraler Normalisierung $u(w)$ (Abschnitt~\ref{sec:phi}) bilde
\[
\Phpref(w)=\bigl(z_0+\rho(w)\,u(w),\ |w|\bigr),
\qquad
\rho(w)=2^{-|w|},\quad z_0=-1-\mathrm{i}/12.
\]
Vergleiche die Verteilung in $\mathcal{M}=\C\times\R$ für
\begin{enumerate}[nosep]
  \item reale Primvierlings-EABC-Wörter $w(p)$ aus $Q(p)=(p,p+2,p+6,p+8)$,
  \item Zufallswörter gleicher Längenverteilung aus $\SigmaEABC^{|w(p)|}$.
\end{enumerate}
\end{definition}

\begin{bemerkung}[Erfolgskriterium]
\emph{Erfolg} wäre eine sichtbare Trennung in zwei chirale Röhren
$T_{\mathrm{ABCE}}$ und $T_{\mathrm{CEAB}}$ (Phasen $+\mathrm{i}$ bzw.\
$-\mathrm{i}$ nach Normalisierung) gegenüber diffuser Zufallsdaten.
\emph{Kein Erfolg} wäre Überlappung oder bloße Trivialität
(alle Vierlinge tragen nur die Wörter $\mathrm{ABCE}$/$\mathrm{CEAB}$).
\end{bemerkung}

\begin{remark}[Nichtbehauptung]
Der Skripttest \texttt{collatz\_phi\_pref\_test.py} ist ein
\emph{explorativer} Datentest. Ein positives Clustering-Bild stützt
höchstens die geometrische Lesart von $\Phpref$ auf EABC-Wörtern;
es beweist weder $\Phi=\Phpref\circ\kappa$ noch $\Einf=\emptyset$.
\end{remark}

%======================================================================
\section{Abgrenzung}
%======================================================================

\begin{center}
  \begin{tabular}{@{}lll@{}}
    \toprule
    Objekt & Status & Referenz \\
    \midrule
    Drei Ebenen (algebraisch/geometrisch/dynamisch) & Methodik & dieses Notizblatt \\
    $z_0=\zeta(-1)(12+\mathrm{i})$, $K_{12}$ & exakt & \texttt{Staudt.tex} \\
    ABCE/CEAB mod $12$ & bewiesen & \texttt{Projektionszeuge\_Primvierling.tex} \\
    $\overline{\mathrm{ABCE}}=\mathrm{CEAB}$, $\chi_H=\pm\mathrm{i}$ & numerisch & \texttt{energiedoku\_eabc\_c4\_kohaerenz.tex} \\
    Primvierlingsellipse $(a,b,e)=(4,2,\sqrt{3}/2)$ & geometrische Kodierung & Ebene~B$'$ \\
    Globale ABCE/CEAB-Zählasymmetrie & offen, kein Bias bis $10^8$ & \texttt{Projektionszeuge\_Primvierling.tex} \\
    Bernoulli-Normschale $\mathcal{I}(n)$ & No-Go (Lyapunov) & \texttt{collatz\_schlussartikel\_arxiv.tex} \\
    Zwölf-Linie $\mathcal{Z}_{12}$ & exakt charakterisiert & \texttt{Staudt.tex} \\
    Stirling$\to|B_{2n}|$ & Mathlib-Baustein, Kombination offen & \texttt{collatz\_lean\_mathlib\_update.tex} \\
    Stereographie $\sigma:S^2\to\Chat$ & Standard & \texttt{Edinburg.py} \\
    $\Einf=\emptyset$ $\Leftrightarrow$ Collatz & präzise, offen & \texttt{collatz\_equivalenz\_e\_infty.tex} \\
    $\Phpref:\SigmaEABC^{<\omega}\to\mathcal{M}$ & wohldefiniert (Wörter) & dieses Notizblatt, \texttt{PrefProjection.lean} \\
    Bernoulli-Uhr $U(m)$ & Gedankenmodell & Abschnitt~\ref{sec:bernoulli-uhr}, \texttt{BernoulliClock.lean} \\
    $\Phi=\Phpref\circ\kappa:\Ztwo\to\mathcal{M}$, $\mathcal{E}_\infty\subset\mathcal{M}$ & Spekulation & dieses Notizblatt \\
  \bottomrule
  \end{tabular}
\end{center}

\medskip
\noindent
Dieses Gedankenexperiment verbindet den Kepler-Ellipsen-Anker
$e_{\mathrm{PV}}=\sqrt{3}/2$, den $\zeta(-1)$-Punkt $z_0$ und die
Riemann-Kompaktkugel zu einer gemeinsamen geometrischen Sprache.
Es ersetzt weder die Uniformitätslücke
(\texttt{collatz\_offene\_punkte.md}, §2) noch die Brücke
\[
\text{endliche EABC-Grammatik}\Rightarrow\Einf=\emptyset.
\]
Collatz bleibt offen.

\end{document}
